\chapter{Conclusion}
\label{ch:conclusion}

This thesis compared VQGAN and LlamaGen VQ-16 to determine how two pretrained
image tokenizers with the same $16\times16$ output grid differ beyond
reconstruction quality. The study measured reconstruction, empirical codebook
usage, token stability under global noise, encoder and decoder locality, and
the response to distribution shift. Each analysis isolated the tokenizer from
the autoregressive model that would consume its tokens.

The reconstruction analysis established the clean baseline and reproduced the
expected advantage of LlamaGen. On ImageNet validation, LlamaGen achieved
$20.79$ rather than $19.65$ PSNR, $0.558$ rather than $0.489$ SSIM, $0.228$
rather than $0.286$ LPIPS, and $2.19$ rather than $5.06$ FID. The codebook
measurements revealed a larger difference between the checkpoints. LlamaGen
used all 16{,}384 entries observed in this evaluation, while VQGAN used about
972. The nominal vocabulary size therefore did not describe the vocabulary
used by each checkpoint.

The analysis of global image perturbations found that neither tokenizer
preserved exact token identity under modest Gaussian noise. At $\sigma=0.1$,
79.5\% of LlamaGen positions and 75.1\% of VQGAN positions differed from their
clean assignments. The decoded images still retained their broad visual
structure, and LlamaGen kept its reconstruction advantage at every tested
noise level. Token assignments can therefore change much faster than the
reconstructed image. Noise also concentrated code usage in both models without
greatly changing the raw number of active codes.

The analysis of local image perturbations found that the encoder response was
not confined to the perturbed patch. At low patch noise, the far-field token
flip rate was 17.4\% for VQGAN and 17.8\% for LlamaGen. Masking one image patch
changed 56.3\% of VQGAN positions and 49.5\% of LlamaGen positions outside the
corresponding token region. A local change in image space can therefore alter
the discrete representation well beyond that region.

The decoder analysis found a more concentrated response. Directly editing a
local token patch produced the largest pixel changes near the edited region,
although measurable changes remained outside it. The same pattern held on the
slightly out-of-distribution ImageNet-V2 dataset, where the largest change in
patch LPIPS relative to ImageNet validation was 0.0068. Locality is therefore
preserved approximately rather than exactly when the decoder maps token edits
back to pixels.

When tested on out-of-distribution datasets, both VQGAN and LlamaGen retained
almost all codes that were active on ImageNet. They did not switch to a
substantially different set of active tokens. Instead, they changed how often
they used those tokens. ImageNet-V2 and ObjectNet remained close to the
ImageNet token distribution, while sketches, medical images, and document
pages produced larger changes. RVL-CDIP produced the largest
Jensen--Shannon divergence for both tokenizers. Even these large changes in
token frequency did not consistently produce large losses in paired
reconstruction quality.

Overall, better reconstruction and broader codebook usage did not imply greater
stability under noise. LlamaGen reconstructed images more accurately and used
far more codes, yet its token assignments changed at a rate similar to VQGAN's.
The distribution-shift results point to the same distinction. Large changes in
the token stream can coexist with modest changes in reconstruction. A
tokenizer therefore determines more than how faithfully an image is rebuilt;
it also determines the sequence distribution presented to the autoregressive
model.

\section{Limitations and Future Work}
\label{sec:limitations_future_work}

This study compared one pretrained VQGAN checkpoint with one pretrained
LlamaGen VQ-16 checkpoint. Because the models differ in architecture, training
procedure, objective, codebook dimensionality, and normalization, the observed
differences describe the two complete systems. They do not identify which
design choice caused each result. A direct mechanistic study should modify one
tokenizer property at a time, retrain the affected models, and repeat the same
measurements. This would separate the effects of codebook size, training
objective, and architecture.

Both tokenizers also require a fixed $256\times256$ input. Resizing and cropping
can alter datasets whose native resolution, aspect ratio, or layout differs
from ImageNet, especially document and medical images. Future comparisons
should include tokenizers with adaptive input resolution and test whether the
same distribution-shift patterns remain without a fixed-size preprocessing
step.

The analyses of global and local perturbations used the same seeds and
intervention levels for both models, but they did not save a shared manifest of
the exact noise arrays and patch locations. This avoidable limitation was
identified after data collection. A replication should export every
intervention once and apply the same saved intervention to both tokenizers.

Finally, the study did not train or evaluate an autoregressive model on the
measured token streams. It therefore cannot show how token instability,
concentrated code usage, or distribution-dependent token frequencies affect
generation. Future work should train or evaluate autoregressive priors on these
streams and measure generation quality and robustness directly. The spatial
analysis should also be repeated with matched dataset sizes, since unequal
sample counts affect the current positional-entropy estimates.
